\section{Implementative challenges encountered}

The most challenging aspect of the implementation was the development and the results enhancing for the \textbf{Key Scan} physical operator, in particular we obtained very poor performance in \texttt{GALOIS F} and \texttt{GALOIS WO} 
logical plans.
This is mainly due to the fact in the concept of Key Scan we first ask to the LLM to retrieve data using only the primary key for each table, so fo instance if we have the \texttt{venezuela\_presidents} table with \texttt{president\_name} defined as table's Primary Key, in the first phase we will ask to the LLM to retrieve the presidents of all word and only in the second phase we will apply the original query to the result set returned by the LLM;
but in principle if the LLM didn't return any president from Venezuela the final result set post query filtering will be empty.
This behavior is common to all datasets used in the \texttt{GALOIS F} and \texttt{GALOIS WO} versions, so we had to find a way to improve the Key Scan's  retrieval process effectiveness.

Another issue we faced was managing DuckDB connections during execution, since several classes open connections to retrieve column data types or, more broadly, the database schema.
For address this problem, the system simply closes the connection to the database when the actions for that specific connection are completed, thus avoiding having too many open connections at the same time.

!!!!NOMINARE IL FATTO DI POTER OTTIMIZZARE LA STRUTTURA E LA SEMANTICA DEI PROMPT, VISTO CHE QUELLI USATO DA GALOIS NON SONO EFFICACI !!!!
